\chapter{总结与展望}

\section{工作总结}

本作品围绕“基于多类格困难假设的抗量子身份基匹配加密通用设计和优化实现”完成了从理论构造、
密码原语到应用入口的完整实现。系统以两层匿名HIBE承载接收身份和认可发送身份，以IBS认证
发送者，并通过短种子混合加密承载不同长度的应用消息。完整密文进一步绑定一次性验签公钥和
全部传输字节，使身份条件、载荷与外层字节认证处于同一处理流程。

在算法层，作品实现LATTE、Lai两种HIBE，以及BNN＋ML-DSA、BNN＋Falcon、Luo、SRGG、
Sageloli五种常规IBS，形成十种端到端组合；在此基础上又接入Wang--PW目标参数矩阵SIS后端。
三组代表分别服务于交互式定向消息、后台异步投递和固定机构间周期报送。所有组合复用相同的公共
五算法、对象边界、状态管理、命令行和浏览器入口。

在优化层，作品针对完整运行中出现的主要瓶颈提出三项实现方法：Sageloli结构化采样保留环结构，
把大规模相关扰动转化为逐频率低秩计算；SRGG范围感知编码按字段实际取值范围压缩四类对象；
Lai接收会话复用固定身份对的叶密钥，将重复委派移到准备阶段。

在软件层，后端通过稳定接口、能力声明和规范编解码接入。组合预检在建立系统前确认身份宽度、
签名输入、载荷容量、对象格式和运行依赖；固定发钥服务与状态文件保存支持跨进程恢复；CLI和
Web界面完成密钥建立、角色导出、定向加密、匹配解密及错误输入拒绝。

\section{实验结论}

十种常规组合在Ultra~9~275HX统一环境下全部完成，共通过20次加密、30次解密、10次新进程
恢复和20次外层拒绝。A组重复加密、解密分别为13.5毫秒和59.4毫秒；B组分别为11.3秒和
68.7毫秒，额外时间集中在发送端SRGG签名；C组分别为89.4秒和409秒。实验条件见第七章。

横向比较进一步表明，HIBE主要决定接收状态、种子密文长度和解密委派，IBS主要决定发送方密钥生成、
签名生成和发送方密钥尺寸。LATTE五组首次解密为95.1至616毫秒，Lai五组为412至440秒；LATTE
接收状态约125~KiB，Lai接收状态约19.5~MiB。固定IBS后，由LATTE换为Lai使完整密文增加约
1.137~MiB；Sageloli组合的完整密文进一步达到3.21和4.34~MiB。

C组阶段分解实验中，Lai委派约占完整解密的92.5\%。接收会话把该阶段前移后，会话准备约510秒，
后续解密中位数为28.6秒，提速约18.4倍。Sageloli PARAM-I、II均完成目标参数发钥、签名和公开验证，结构化
实现的峰值RSS分别约1.36~GiB和6.19~GiB。三组SRGG参数机构状态的紧凑长度为JSON的18.3\%至
19.8\%。Wang--PW完成目标规模建立、身份发钥、独立签名、公开验签及LATTE组合，其大对象和
服务器级资源需求也说明通用接口需要同时容纳轻量与重型后端。

密文组成分析表明，A组主要开销来自Lamport一次性认证，C组主要开销来自Sageloli签名和Lai
种子密文。优化单个对象能否显著改变完整传输量，取决于该对象在整包中的占比。身份错配、字节篡改、
截断、追加、内部签名错误和跨进程恢复测试则覆盖了统一协议的主要功能路径。

\section{后续工作}

后续研究将继续提高通信效率。围绕短消息中的固定认证开销和大签名传输，探索更紧凑的密文
组织与认证方法，降低消息封装、存储和传输所需的字节数，使身份匹配加密适应更广泛的网络
带宽与终端资源条件。

算法体系将向更多格结构和应用需求扩展。结合不同HIBE的身份层级与委派特点，以及不同IBS的
签名长度、密钥规模和运算方式，形成更灵活的组合选择；进一步研究计算复用、批量任务与
并行处理，使算法特征与消息规模、通信频率相适应。

面向长期应用，将深化身份与密钥管理，研究有效期、密钥更新与撤销如何同双方身份条件配合，
并推动系统与组织目录、设备管理和消息服务集成。通过完善机构和用户之间的密钥分发与通信
管理，支持更复杂的组织协作和持续运行。

系统将进一步发展为支持多类格算法、适应不同资源与通信条件的身份基匹配加密平台，为机构
定向消息、设备管理和跨域数据交换提供统一的密码服务。
